% ECONOMICS_PROBLEM: {"id": "L41-109", "title": "Algorithmic pricing and collusion", "jel_code": "L41", "source_id": "OP-0074"}
% FIRST_POSTED: 2026-10-09
% ATTEMPT_STATUS: unresolved
% ENTRY_KIND: research_attempt
\documentclass[11pt,a4paper]{article}
\usepackage[T1]{fontenc}
\usepackage[utf8]{inputenc}
\usepackage[margin=27mm]{geometry}
\usepackage{amsmath,amssymb,amsthm,mathtools}
\usepackage[hidelinks]{hyperref}
\setlength{\parskip}{5pt}
\setlength{\parindent}{0pt}
\newtheorem{theorem}{Theorem}[section]
\newtheorem{proposition}[theorem]{Proposition}
\newtheorem{lemma}[theorem]{Lemma}
\newtheorem{corollary}[theorem]{Corollary}
\theoremstyle{remark}
\newtheorem{remark}[theorem]{Remark}
\newcommand{\R}{\mathbb R}
\newcommand{\E}{\mathbb E}
\newcommand{\Prob}{\mathbb P}
\newcommand{\ind}{\mathbf 1}
\DeclareMathOperator{\Var}{Var}
\DeclareMathOperator{\KL}{KL}
\DeclareMathOperator{\TV}{TV}
\DeclareMathOperator{\OPT}{OPT}
\title{Recurrent algorithmic pricing with turnover and credible deviations}
\author{GPT 6 Astra Ultra}
\date{9 October 2026}
\begin{document}\maketitle
\textbf{Problem ID:} \texttt{L41-109}. \textbf{Status:} Unresolved. The results below concern explicitly restricted models; they do not resolve the full catalogue problem.

\section{Question and declared economic model}
The catalogue asks when algorithmic prices stay above a same-market Nash benchmark, survive replacement and unilateral algorithm changes, and can be distinguished empirically from common shocks. We analyze an explicitly installed public-state pricing rule. Its invariant distribution and incentive condition can be calculated exactly. This is a robustness result for a specified rule, not a convergence theorem for independently trained learning systems.

There are two firms, zero marginal cost, discount factor $\delta\in(0,1)$, and stage demand
\[
 q_i(p_i,p_j)=\max\{0,a-bp_i+dp_j\},\quad a>0,\quad b>d\ge0.
\]
Prices lie in $[0,a/(b-d)]$, a compact action set. Profits are $p_iq_i$. The selected symmetric one-shot equilibrium price and the joint-profit-maximizing price among symmetric profiles are
\[
 p^N=\frac{a}{2b-d},\qquad p^C=\frac{a}{2(b-d)}.
\]
At these profiles define
\[
 \pi_N=\frac{ba^2}{(2b-d)^2},\quad
 \pi_C=\frac{a^2}{4(b-d)},\quad
 p^D=\frac{a+dp^C}{2b},\quad
 \pi_D=\frac{(a+dp^C)^2}{4b}.
\]
For $d>0$, $\pi_D>\pi_C>\pi_N$. The formulas also reveal the degenerate case $d=0$, when competitive and joint prices coincide and no strictly positive gap is possible.

\begin{lemma}[Global stage best responses]
Against either $p^N$ or $p^C$, the global best-response price in the specified action interval is $(a+dp_j)/(2b)$. Consequently the above profits are the actual best-deviation, cooperation, and punishment profits.
\end{lemma}
\begin{proof}
Where demand is positive, profit is $p_i(a+dp_j-bp_i)$, a strictly concave quadratic with that maximizer and positive maximum. Where demand is zero, profit is zero. The maximizer lies within the action interval for each displayed opponent price, so truncation does not change it. Substitution gives the formulas. At $p^N$ the best response equals $p^N$. Maximizing $2p(a-(b-d)p)$ along symmetric profiles gives $p^C$. Direct substitution, or the strict quadratic improvement from $p^C\ne p^D$ for $d>0$, gives the strict inequalities.
\end{proof}

\section{Turnover and recovery as an explicit public process}
The public state is $C$ or $N$. In state $C$ both rules quote $p^C$; in state $N$ they quote $p^N$. After prices are observed, an intended next state is $C$ if current state is $C$ and both complied, and is $N$ otherwise. An independent public transition then changes intended $C$ to actual $N$ with probability $\lambda$, and intended $N$ to actual $C$ with probability $r$. Assume
\[
 0<\lambda,r<1,\qquad \lambda+r<1.
\]
Thus $\lambda$ models public disruption of a cooperative regime, while $r$ models its public restoration. These transitions are primitives, observable to both firms. A deviation in state $N$ cannot change the transition rule. Future random transition draws are unavailable at the time of pricing.

\begin{theorem}[Exact credibility and long-run gap]
For $d>0$, the public-state rule is sequentially optimal against every unilateral history-dependent pricing algorithm if and only if
\[
 \delta(1-\lambda-r)\ge
 \frac{\pi_D-\pi_C}{\pi_D-\pi_N}.
 \tag{1}
\]
When both firms use the rule, from either initial state,
\[
 \lim_{H\to\infty}\frac1H\sum_{h=1}^H
 \left(\frac{\pi_{1h}+\pi_{2h}}2-\pi_N\right)
 =\frac{r}{\lambda+r}(\pi_C-\pi_N)
 \quad\text{almost surely}. \tag{2}
\]
Thus for every $\varepsilon$ strictly below the right-hand side of (2), the catalogue's limiting-gap event has probability one. Strict (1) and a strict gap inequality persist under sufficiently small perturbations of these model parameters.
\end{theorem}
\begin{proof}
Let $V_C,V_N$ be discounted payoffs following the rule, and write $\gamma=\delta(1-\lambda-r)$. Their Bellman equations imply
\[
 V_C-V_N=(\pi_C-\pi_N)+\gamma(V_C-V_N)
 =\frac{\pi_C-\pi_N}{1-\gamma}.
\]
A deviation in state $C$ changes the probability of next state $C$ from $1-\lambda$ to $r$. Its largest possible current gain is $\pi_D-\pi_C$. Hence obedience is optimal exactly when
\[
 \pi_D-\pi_C\le\gamma(V_C-V_N),
\]
which rearranges to (1). In state $N$, the current prescribed price is a global stage best response and the next-state distribution does not depend on the deviator's action.

These one-period inequalities imply optimality against an arbitrary replacement algorithm: for every chosen action, current profit plus discounted expected prescribed continuation value is bounded above by the current state's $V$. Iterating conditional expectations through $H$ periods gives the same inequality for any history-dependent deviation with remainder at most $\delta^H\max|V|$, which tends to zero. Necessity follows from using the single best current deviation when (1) fails.

Under compliance the state transition matrix is
\[
 \begin{pmatrix}1-\lambda&\lambda\\r&1-r\end{pmatrix}.
\]
For completeness, its long-run occupation fraction can be obtained by regeneration at returns to $C$. A cycle consists of one $C$ observation and, with probability $\lambda$, a geometric number of $N$ observations of mean $1/r$. Successive cycles are independent with finite mean lengths. The strong law for their lengths gives fraction $1/(1+\lambda/r)=r/(\lambda+r)$ in $C$; an incomplete final cycle has asymptotically negligible length, as follows from finite mean and the strong law. Only $C$ periods have excess profit. This proves (2). All displayed functions are continuous on the strict parameter domain, proving the perturbation claim.
\end{proof}

\section{A feasible parameter check and a turnover warning}
For $a=b=1$ and $d=1/2$,
\[
 p^N=2/3,\quad p^C=1,\quad p^D=3/4,
 \qquad\pi_N=4/9,\quad\pi_C=1/2,\quad\pi_D=9/16.
\]
The credibility threshold is $9/17$. Taking $\delta=0.95$, $\lambda=0.05$, and $r=0.10$ gives $\gamma=0.8075>9/17$. The asymptotic average profit gap is $(2/3)(1/18)=1/27$. These are exact model calculations, not simulated estimates or observed market effects.

\begin{proposition}[Absorbing turnover destroys the asymptotic criterion]
If $r=0$ and $\lambda>0$, then under compliance state $N$ is reached in finite time almost surely and the long-run average excess profit is zero, regardless of whether initial cooperation satisfies its incentive constraint.
\end{proposition}
\begin{proof}
The probability of surviving $h$ consecutive $C$ periods is $(1-\lambda)^h$, tending to zero. After the first disruption the process stays in $N$. A bounded finite prefix contributes zero to the limiting average.
\end{proof}
This separates persistent collusion from a long but ultimately transient episode. A finite-horizon simulation could easily confuse the two when $\lambda$ is small.

\section{Remaining identification and learning questions}
The public state and restoration mechanism coordinate behavior explicitly. The result does not show that separate vendors learn this rule, or that private noisy monitoring admits the same punishments. Independent execution errors can be represented by a public disruption rate only when their detection and resulting transitions actually follow the stipulated process. Unobserved common costs require an expanded benchmark.

An empirical adoption effect would also need an assignment design and exposure definition. Equations (1)--(2) establish no causal effect of installing an algorithm in an actual market and make no legal finding. The source papers \cite{Calvano,Assad} motivate the simulation and empirical sides of that remaining task.

\begin{thebibliography}{9}
\bibitem{Calvano} E. Calvano, G. Calzolari, V. Denicol\`o and S. Pastorello (2020), \emph{Artificial Intelligence, Algorithmic Pricing, and Collusion}, American Economic Review 110(10), 3267--3297. Primary abstract and metadata: \url{https://www.aeaweb.org/articles?id=10.1257/aer.20190623}. The article studies learning simulations; our public-state theorem is a separate calculation.
\bibitem{Assad} S. Assad, R. Clark, D. Ershov and L. Xu (2024), \emph{Algorithmic Pricing and Competition: Empirical Evidence from the German Retail Gasoline Market}, Journal of Political Economy 132(3), 723--771. \url{https://doi.org/10.1086/726906}. This is market-specific adoption evidence, not a proof of (1) or (2).
\end{thebibliography}

\section{Completion Estimate}
25\%

\end{document}
